\documentclass[10pt,a4paper]{amsart}
\usepackage{amsmath,amssymb,amsthm,graphicx}
\usepackage{verbatim}
\newenvironment{reference}{\comment}{\endcomment}
%\newenvironment{reference}{}{}
\newenvironment{slogan}{\comment}{\endcomment}
\newenvironment{history}{\comment}{\endcomment}

% For commutative diagrams we use Xy-pic
\usepackage[all]{xy}

% We use 2cell for 2-commutative diagrams.
\xyoption{2cell}
% Theorem environments.
%
\theoremstyle{plain}
\newtheorem{theorem}[subsection]{Theorem}
\newtheorem{proposition}[subsection]{Proposition}
\newtheorem{lemma}[subsection]{Lemma}

\theoremstyle{definition}
\newtheorem{definition}[subsection]{Definition}
\newtheorem{example}[subsection]{Example}
\newtheorem{exercise}[subsection]{Exercise}
\newtheorem{situation}[subsection]{Situation}

\theoremstyle{remark}
\newtheorem{remark}[subsection]{Remark}
\newtheorem{remarks}[subsection]{Remarks}

\numberwithin{equation}{subsection}

% Macros
%
\def\lim{\mathop{\mathrm{lim}}\nolimits}
\def\colim{\mathop{\mathrm{colim}}\nolimits}
\def\Spec{\mathop{\mathrm{Spec}}}
\def\Hom{\mathop{\mathrm{Hom}}\nolimits}
\def\Ext{\mathop{\mathrm{Ext}}\nolimits}
\def\SheafHom{\mathop{\mathcal{H}\!\mathit{om}}\nolimits}
\def\SheafExt{\mathop{\mathcal{E}\!\mathit{xt}}\nolimits}
\def\Sch{\mathit{Sch}}
\def\Mor{\mathop{\mathrm{Mor}}\nolimits}
\def\Ob{\mathop{\mathrm{Ob}}\nolimits}
\def\Sh{\mathop{\mathit{Sh}}\nolimits}
\def\NL{\mathop{N\!L}\nolimits}
\def\CH{\mathop{\mathrm{CH}}\nolimits}
\def\proetale{{pro\text{-}\acute{e}tale}}
\def\etale{{\acute{e}tale}}
\def\QCoh{\mathit{QCoh}}
\def\Ker{\mathop{\mathrm{Ker}}}
\def\Im{\mathop{\mathrm{Im}}}
\def\Coker{\mathop{\mathrm{Coker}}}
\def\Coim{\mathop{\mathrm{Coim}}}

% Boxtimes
%
\DeclareMathSymbol{\boxtimes}{\mathbin}{AMSa}{"02}

%
% Macros for moduli stacks/spaces
%
\def\QCohstack{\mathcal{QC}\!\mathit{oh}}
\def\Cohstack{\mathcal{C}\!\mathit{oh}}
\def\Spacesstack{\mathcal{S}\!\mathit{paces}}
\def\Quotfunctor{\mathrm{Quot}}
\def\Hilbfunctor{\mathrm{Hilb}}
\def\Curvesstack{\mathcal{C}\!\mathit{urves}}
\def\Polarizedstack{\mathcal{P}\!\mathit{olarized}}
\def\Complexesstack{\mathcal{C}\!\mathit{omplexes}}
% \Pic is the operator that assigns to X its picard group, usage \Pic(X)
% \Picardstack_{X/B} denotes the Picard stack of X over B
% \Picardfunctor_{X/B} denotes the Picard functor of X over B
\def\Pic{\mathop{\mathrm{Pic}}\nolimits}
\def\Picardstack{\mathcal{P}\!\mathit{ic}}
\def\Picardfunctor{\mathrm{Pic}}
\def\Deformationcategory{\mathcal{D}\!\mathit{ef}}
\title{The valuative criterion for algebraic spaces is preserved by composition}
\author{The Stacks Project Authors}
\date{stacks-project:spaces-morphisms.tex:local-snapshot-20261012; GNU FDL 1.2 or later}
\begin{document}
\maketitle
\setcounter{subsection}{1}
\noindent\textit{Editorial note.} All delivered mathematical text is the author's own, copied from the pinned Stacks Project chapter file spaces-morphisms.tex; no statement or proof step was rewritten, reordered or removed.\par\medskip
\noindent\textit{Editorial note.} The existence case is presented through the displayed two-step diagram and then spelled out in words; both are inside the counted proof span.\par\medskip
\noindent\textit{Editorial note.} The counted spans carry no references at all, so the card needs no navigation stubs beyond the definition pointer implied by the statement shape.\par\medskip
\noindent\textit{Editorial note.} Preamble: the three native excerpt spans of sources/stacks-spaces-morphisms/preamble.tex (lines 8-13, 14-18 and 154-217). The chapter file itself carries no preamble text.\par\medskip
\noindent\textit{Allowed standard prerequisites (editorial).}\begin{itemize}
\item The shape of the valuative criterion diagram. Definition definition-valuative-criterion, used at proof lines 8738-8745. Fixes the solid diagram over the composition Z <- Y <- X.
\item Uniqueness and existence parts of the criterion for the two factors. used at proof lines 8746-8750 and 8765-8768. Are applied successively to build and compare the dotted arrows.
\item Valuation rings under field extension and domination are again of the same kind. used at proof lines 8754-8762. Allows the second application of the criterion to use the field extension produced by the first.
\end{itemize}\medskip
\begin{lemma}
\label{lemma-composition-valuative-criteria}
The composition of two morphisms of algebraic spaces which satisfy the
(existence part of, resp.\ uniqueness part of) the valuative criterion
satisfies the (existence part of, resp.\ uniqueness part of) the valuative
criterion.
\end{lemma}

\begin{proof}
Let $f : X \to Y$, $g : Y \to Z$ be morphisms of algebraic spaces over the
scheme $S$. Consider a solid commutative diagram of the following shape
$$
\xymatrix{
\Spec(K) \ar[dd] \ar[r] & X \ar[d]^f \\
& Y \ar[d]^g \\
\Spec(A) \ar[r] \ar@{-->}[ru] \ar@{-->}[ruu] & Z
}
$$
If we have the uniqueness part for $g$, then there exists at
most one north-west dotted arrow making the diagram commute.
If we also have the uniqueness part for $f$, then we have
at most one north-north-west dotted arrow making the diagram
commute. The proof in the existence case comes from contemplating
the following diagram
$$
\xymatrix{
\Spec(K'') \ar[r] \ar[dd] &
\Spec(K') \ar[r] &
\Spec(K) \ar[r] &
X \ar[d]^f \\
& & & Y \ar[d]^g \\
\Spec(A'') \ar[r] \ar[rrruu] &
\Spec(A') \ar[r] \ar[rru] &
\Spec(A) \ar[r] &
Z
}
$$
Namely, the existence part for $g$ gives us the extension $K'$, the
valuation ring $A'$ and the arrow $\Spec(A') \to Y$, whereupon
the existence part for $f$ gives us the extension $K''$, the
valuation ring $A''$ and the arrow $\Spec(A'') \to X$.
\end{proof}

\end{document}
